% SPDX-License-Identifier: Apache-2.0
% covers: FlashWin
\datasheet{FlashWin}{A read-only window onto an SPI flash: a read on the bus becomes a fast read on the wires, and the word comes back.}

\dsfacts{\code{lib/parts/src/flashwin.rs}}{\code{txhdl\_parts::flashwin::FlashWin}}%
{\code{//docs:examples}}

\subsection*{Function}

The SPI master in \code{spi} is what a program drives a byte at a
time. \code{FlashWin} is the other way of reaching the same chip: an
address range in which an ordinary read is answered with what the
flash holds there, with no registers to drive and nothing for software
to sequence.

That is what a bootloader needs, and what a core fetching from the bus
needs, because neither can run software to read the software it is
about to run.

A read at an address sends the fast read command \code{0x0b}, the
three bytes of that address with its low two bits cleared, one byte
the chip throws away, and then takes four, which are the word, lowest
address in the lowest byte. Nine bytes on the wires for four on the
bus.

\subsection*{Parameters}

\begin{center}\footnotesize
\begin{tabular}{@{}lp{0.7\textwidth}@{}}
\toprule
Parameter & Meaning \\
\midrule
\code{DIV} & The clock: a half of a bit takes \code{DIV + 1} cycles,
so the flash is clocked at the system's rate over
\code{2 * (DIV + 1)}. The tables use 1, as the example does; a board
divides further. \\
\code{I} & The width of the AXI identifier on \code{req}, \code{ans}
and \code{rb}, in bits. The tables use 2, as every design here does. \\
\bottomrule
\end{tabular}
\end{center}

\subsection*{Register map}

None. The window has no registers: every address in its range is a
word of the flash, and the range's base is the router's business
rather than the part's, as it is for the data memory.

\dstables{FlashWin}

\subsection*{Behaviour}

\begin{itemize}
\item A read is taken when nothing is on the wires and no answer is
waiting. Its address is kept with the low two bits cleared, since a
word is what the bus asks for.
\item The nine bytes go out under one select. The select is low from
the cycle the read is taken and rises when the ninth byte is done, so
each read is a command the chip sees whole, which is what a real chip
requires.
\item The wires are mode 0, and only mode 0: the clock idles low, the
chip takes a bit as the clock rises and the master moves one as it
falls. Every flash answers a fast read in that mode, so the part needs
neither the master's two mode bits nor a register to hold them.
\item The four bytes of data are assembled as they arrive, the first
into the low eight bits. The answer goes out on \code{rb} in the first
cycle the channel has room after the ninth byte.
\item A write is taken, its beat is dropped and it is answered
\code{SLVERR}. The window is read only, and saying so is what makes it
safe to fetch from: nothing on the bus can change what the instruction
stream is reading, so the question of what a fetch sees when a write
lands in the same cycle does not arise. That question is issue 268.
\item \code{rst} stops whatever is on the wires and clears the answer
and the held write.
\end{itemize}

\subsection*{Verification}

\begin{itemize}
\item \code{//lib/examples:ex\_flashwin} puts the window behind an AXI
link, reads two words the chip model holds and checks each against
what was put there, then writes and finds \code{SLVERR}.
\item The chip is \code{FlashDevice} from \code{spi}, the model that
already checks the master, so the two ways of reaching a flash are
checked against one description of one chip rather than two that could
drift apart.
\item The build simulates the netlist against that run under nvc and
Verilator: \code{//docs:flashwin\_sim\_flashwin\_tb\_test} and
\code{//docs:flashwin\_vsim\_test}.
\end{itemize}

\subsection*{Used by}

Nothing in the tree yet. It is meant for the board's configuration
flash, which issue 312 puts on the board through \code{STARTUPE2}, and
for the bootloader of issue 135.

\subsection*{Limits}

A word costs nine bytes on the wires, which is about three hundred
cycles at the divider the example uses: slow beside a memory and fast
beside a resynthesis. There is no cache in front of it and no burst:
every word is a command of its own, and a read of several beats is not
served. Only mode 0, and only the fast read command, so a chip that
wants another command for another purpose wants the master instead.
The address is 24 bits, which is the three the command carries, so the
window reaches 16\,MiB. Nothing is refused by address: every address
in the range is asked of the chip, and a chip that has nothing there
answers what it answers.
